% ---------------------------------------------------------------------------
% Preamble for building this book with the tufte-book LaTeX class.
% Loaded via `include-in-header` in _quarto.yml.
% ---------------------------------------------------------------------------

% --- Fix: tufte-latex vs the modern LaTeX kernel ---------------------------
%
% tufte-common.def defines:
%
%   \newcommand{\allcaps}[1]{\allcapsspacing{\MakeTextUppercase{#1}}}
%   \newcommand{\smallcaps}[1]{\smallcapsspacing{\MakeTextLowercase{#1}}}
%
% Two problems with current TeX Live:
%
%   1. \MakeTextUppercase / \MakeTextLowercase come from `textcase`, which
%      clashes with the kernel's own case-changing machinery.
%   2. Both are declared with \newcommand, so they are *short* commands and
%      reject any argument containing a paragraph break. The title page and
%      the running heads both trip this, giving:
%        "Argument of \MakeTextUppercase has an extra }."
%
% tufte redefines these again during font setup, so patching them in the
% preamble is not enough — it must happen at \begin{document}.
% The whole chain is short: \allcaps -> \allcapsspacing -> \MakeTextUppercase,
% and the kernel's \MakeUppercase is short too. So we define replacements that
% hand the argument to no argument-taking command at all: \uppercase is a TeX
% primitive and tolerates \par, and \scshape is a declaration rather than a
% command with an argument. We lose tufte's letterspacing on caps; we gain a
% document that compiles.
\makeatletter
\AtBeginDocument{%
  \long\def\allcaps#1{{\uppercase{#1}}}%
  \long\def\smallcaps#1{{\scshape #1}}%
}
\makeatother

% --- Title page ------------------------------------------------------------
% tufte-book's title page prints author, title, then publisher.
\publisher{r-best-practices.njtierney.com}

% --- Quarto compatibility --------------------------------------------------
% Quarto emits \tightlist; tufte-book does not define it.
\providecommand{\tightlist}{%
  \setlength{\itemsep}{0pt}\setlength{\parskip}{0pt}}

% --- tufte-book refuses to define \subsubsection ---------------------------
%
% tufte-common.def replaces \subsubsection with an error, on Bringhurst's
% argument (_The Elements of Typographic Style_, 4.2.2) that a document
% needing a fourth heading level wants restructuring rather than another size
% of type. Using one aborts the build with:
%
%   Class tufte-book Error: \subsubsection is undefined by this class.
%
% This book reaches four levels in one place: Files / Directories / Paths,
% nested under "Files, directories, paths". Define a quiet italic heading
% rather than lose the PDF. \@startsection is used rather than a plain
% \textbf so that the table of contents and hyperref anchors still work.
%
% If the fourth level ever goes away, this can go with it.
\makeatletter
\renewcommand\subsubsection{%
  \@startsection{subsubsection}%
    {3}%                                 level
    {\z@}%                               indent
    {1.4ex \@plus .4ex \@minus .2ex}%    space before
    {0.6ex \@plus .2ex}%                 space after
    {\normalfont\normalsize\itshape}%    style
}
\makeatother

% --- Fix: Quarto's margin packages vs tufte's own --------------------------
%
% Quarto emits, near the end of its preamble:
%
%   \@ifpackageloaded{sidenotes}{}{\usepackage{sidenotes}}
%   \@ifpackageloaded{marginnote}{}{\usepackage{marginnote}}
%
% Both abort the build on this class, because tufte-common.def has already
% defined \sidenote and \marginnote itself:
%
%   "Command \marginnote already defined."
%   "Command '\sidenote' already defined."
%
% Those two packages give an ordinary class the margin apparatus that
% tufte-book is built around, so here they are redundant - and nothing in
% the book calls either command. Declaring them already loaded makes Quarto
% take the empty branch and leaves tufte's versions in place. This has to
% come before Quarto's block, which it does: include-in-header lands earlier
% in the preamble.
\makeatletter
\@namedef{ver@sidenotes.sty}{9999/01/01}
\@namedef{ver@marginnote.sty}{9999/01/01}
\makeatother

% --- Every image runs the full width of its column ------------------------
%
% Images arrive here three ways, and only the first is already full width:
%
%   1. `width=100%` in the .qmd  -> \includegraphics[width=1\linewidth]
%   2. `width=70%`  in the .qmd  -> \includegraphics[width=0.7\linewidth]
%   3. no width at all           -> \pandocbounded{\includegraphics{...}},
%      which only ever *shrinks* an image to fit, never grows it, so a small
%      PNG stays small.
%
% Rather than edit every image in the source - which would also change the
% HTML, where the 70% sizing is wanted - force the width here. This file is
% only loaded by the PDF format, so HTML is untouched.
%
% The passed options `#1` are kept and our width is appended *after* them:
% graphicx uses the last value given for a key, so this beats an explicit
% `width=0.7\linewidth` while preserving `alt=`, which carries the fig-alt
% text into the tagged PDF.
\usepackage{letltxmacro}
\LetLtxMacro{\rbporigincludegraphics}{\includegraphics}
\renewcommand{\includegraphics}[2][]{%
  \rbporigincludegraphics[#1,width=\linewidth,height=\maxdimen,keepaspectratio]{#2}%
}

% Code blocks are wider than tufte's narrow text column, so allow a little
% overflow into the margin rather than letting lines run off the page.
\usepackage{fvextra}
\DefineVerbatimEnvironment{Highlighting}{Verbatim}{
  breaklines, breaknonspaceingroup, breakanywhere, commandchars=\\\{\}
}

% --- Code blocks run out into the margin -----------------------------------
%
% Quarto sets code as `Shaded` (from `framed`) wrapping `Highlighting`, and
% `Shaded` takes its width from \hsize, which in tufte is the narrow text
% column. That squeezes pipelines onto wrapped lines.
%
% Wrapping Shaded in tufte's `fullwidth` does not work: `fullwidth` is a
% `changepage` list environment, and `framed` cannot survive inside a list.
% It fails with "\begin{snugshade} ended by \end{list}".
%
% So widen the measure rather than moving the block: add tufte's own
% \@tufte@overhang to \hsize before `framed` reads it. `framed` keeps its
% page-breaking, which a \parbox-based box would lose, and long code blocks
% here do break across pages.
%
% The \ifdim guard keeps this to top-level code only: widening code set
% inside a callout would push it out through the side of the box.
%
% The comparison has to be against a saved copy of the body width, not
% against \textwidth. Quarto sets callouts as tcolorbox, and tcolorbox
% re-points \textwidth at the box's inner width along with \linewidth and
% \hsize - so `\ifdim\linewidth=\textwidth` is true inside a callout too,
% and every code block gets widened. \rbpbodywidth is set once, at the top
% level, where it records the real measure.
\makeatletter
\newdimen\rbpbodywidth
\AtBeginDocument{\rbpbodywidth=\textwidth}
\let\rbporigShaded\Shaded
\let\rbporigendShaded\endShaded
\renewenvironment{Shaded}{%
  \begingroup
  \ifdim\linewidth=\rbpbodywidth
    \advance\linewidth by \@tufte@overhang
    \advance\hsize by \@tufte@overhang
    \advance\columnwidth by \@tufte@overhang
  \fi
  \rbporigShaded
}{%
  \rbporigendShaded
  \endgroup
}
\makeatother

% --- Figures span the text column and the margin ---------------------------
%
% A plain `figure` in tufte-book is only as wide as the narrow text column,
% which leaves the diagrams in this book too small to read. tufte-book's
% `figure*` is its fullwidth float: it spans the text column, the gutter and
% the margin column, and still sets the caption in the margin.
%
% Quarto always emits `\begin{figure}`, and there is no format option that
% changes that on this class - `fig-column` only affects the HTML. So alias
% the environment: every figure becomes a fullwidth one. Combined with the
% \includegraphics width above, `\linewidth` inside the float is now the
% full measure, so the image fills it.
\makeatletter
\AtBeginDocument{%
  \expandafter\let\expandafter\figure\csname figure*\endcsname
  \expandafter\let\expandafter\endfigure\csname endfigure*\endcsname
}
\makeatother
